\section{H46: Style is a fingerprint, not a conserved charge}

\textbf{The question:} Each agent writes with a style (formatting, punctuation, pronoun rates) and about some content. Does the style stay fixed through every step change while the content moves? If yes, style belongs to the model weights and can identify the agent.

\textbf{The intuition:} A conserved charge is a quantity that no interaction can change, like electric charge in a collision. H46 asked whether style is such a charge of the model, while content is a state that follows the goal field and the room. The test is like checking a particle's charge before and after each collision. A boundary (a goal switch, a context erasure, a room cut) is a collision. If style is conserved, its jump at a boundary looks like an ordinary day-to-day jitter. The result points to a different picture. A fixed charge of the weights is ``dressed'' by a part that lives in the context window. That part grows as the context fills.

\textbf{What we measured:} Style is 20 text features from H13, with length, code and link shares controlled. Content is the style-residualized embedding. For each agent and boundary, we ranked the boundary jump among the agent's own day-to-day jumps nearby. $T$ is the mean percentile: 0.5 means no move, and 1 means the largest jump. We used 32 non-reserved goal periods and five boundary classes. The null is a randomization test on these ranks, plus boundary-level tests. The scheduler is not relevant: the test uses day-level displacement and no timing; a gap-matched placebo removed the weekend gap. The external drive is the intervention itself (goal switches); topic directions were not residualized out. The shared model prior is the object; content uses style-residualized vectors, and the fingerprint was also tested within Anthropic agents only. Common input is not relevant, because there is no copying claim. A synthetic check came first on the real message schedule. It held nominal size and detected a style shift of half a day-jitter at goal switches in 100\% of runs.

\textbf{What we found:} Style is not conserved. At goal switches style moves and content moves more ($T_s=0.69$ [0.61, 0.77], $T_c=0.86$; 21 of 24 switches). Across forced context erasures, timed by the scaffold, style moves (0.56 [0.51, 0.61]) and content does not (0.52). Style also moves for the \#51 Prankster (percentile 1.00) and for \#12 judges. Style holds only where content also holds (roster, scaffold). But style still identifies agents across a goal switch: accuracy 0.76 for style, 0.51 for content, against 0.15 by chance. Style wins at 22 of 24 switches. Post hoc, style drifts away from the agent's mean as its context fills and returns after an erasure. The claim that stands is ``style drifts with context fill and resets at erasure, yet identifies agents better than content''\cred{0.75}, EU 2.44. The drift-and-reset part of this claim comes from a post hoc analysis. The conservation verdict (refuted) and the fingerprint (passed) were pre-registered.

\begin{figure}[h]
\includegraphics[width=\textwidth]{H46-style-conserved-charge/figures/summary_obs.pdf}
\caption{Style moves at goal switches and erasures, but still identifies agents. (a) Each faint dot is one boundary: style percentile against content percentile, both relative to the agent's own day-to-day jumps. Big symbols are class means, and the gray band is ``no move''. (b) Identity accuracy across the same boundaries, trained before and tested after: style beats content at most goal switches.}
\end{figure}

\textbf{What it means:} For an operator: to attribute a day's messages to a model, use the 20 formatting and punctuation rates, not content embeddings. Across a goal switch they identify the agent at 0.76, against 0.51 for embeddings. The same fingerprint can flag an unannounced model swap. For a monitor builder: compare messages at the same context position and in the same speech act. A context erasure alone shifts style (0.56), and an assigned role such as a prankster or a judge shifts it strongly. For a physicist: treat style as a charge of the weights plus a context-held excitation that resets at erasure.

\textbf{Caveats:} The 20 features mix formatting with topic-adjacent rates; digits, uppercase and colons carry most of the goal-switch shift. A function-word stylometry may be conserved; it was not tested. Speech-act genre is not controlled, so the erasure and judge effects can be genre. The Prankster and NE38 are single agents. The rooms class has four boundaries, and the nudger class has one. The information test is observational and has power only in \#51 and \#38. Content uses one embedding model. No confirmation test on reserved goal periods has run yet.
